

We first formalize the difficulty of direct action-space sampling
by characterizing how the probability mass of feasible trajectories
shrinks with the planning horizon.

\begin{lemma}[Vanishing Feasible Mass in Trajectory Space]
\label{lem:vanishing_mass}
Let $\mathcal{X} = \mathcal{S}^H \times \mathcal{A}^H$ be a bounded trajectory space
with ambient dimension
\[
D = H(\dim \mathcal{S} + \dim \mathcal{A}),
\]
equipped with a reference measure $\mu$.
Let $\mathcal{M}_{\mathrm{traj}} \subset \mathcal{X}$ be a compact subset
with intrinsic dimension $d \ll D$.
Then, for any fixed tolerance $\epsilon \in (0,1)$,
there exists a constant $c > 0$ such that
\begin{equation}
\frac{\mu\big(\mathcal{N}_\epsilon(\mathcal{M}_{\mathrm{traj}})\big)}{\mu(\mathcal{X})}
\le \exp(-cH),
\end{equation}
where $\mathcal{N}_\epsilon(\mathcal{M}_{\mathrm{traj}})$ denotes the $\epsilon$-neighborhood
of $\mathcal{M}_{\mathrm{traj}}$.
\end{lemma}

\begin{proof}
We work in the bounded trajectory space
$\mathcal{X} \subset \mathbb{R}^D$
equipped with the Lebesgue measure $\mu$.
Let $\mathcal{M}_{\mathrm{traj}} \subset \mathcal{X}$ be compact
with intrinsic dimension $d$.

Fix any $\epsilon \in (0,1)$ and define
\[
\mathcal{N}_\epsilon(\mathcal{M}_{\mathrm{traj}})
=
\{x \in \mathcal{X} \mid \inf_{y \in \mathcal{M}_{\mathrm{traj}}}
\|x-y\|_2 \le \epsilon \}.
\]

By standard covering number arguments for compact sets,
$\mathcal{M}_{\mathrm{traj}}$ can be covered by at most
$N_\epsilon \le C_1 \epsilon^{-d}$ Euclidean balls of radius $\epsilon$,
where $C_1 > 0$ is independent of $H$.
Thus,
\[
\mathcal{N}_\epsilon(\mathcal{M}_{\mathrm{traj}})
\subset
\bigcup_{i=1}^{N_\epsilon} B(y_i, 2\epsilon).
\]

Each such ball has $D$-dimensional volume proportional to $\epsilon^D$.
Up to multiplicative factors that grow at most subexponentially with $D$,
we obtain
\[
\mu(\mathcal{N}_\epsilon(\mathcal{M}_{\mathrm{traj}}))
\;\lesssim\;
\epsilon^{D-d}.
\]

Since $\mathcal{X}$ is bounded, $\mu(\mathcal{X}) < \infty$.
Moreover, the ambient dimension grows linearly with the horizon,
\[
D = H(\dim \mathcal{S} + \dim \mathcal{A}),
\]
while the intrinsic dimension of feasible trajectories is governed by
physical and semantic constraints.
We assume that
\[
d \le \lambda H
\quad \text{for some } \lambda < \dim \mathcal{S} + \dim \mathcal{A},
\]
so that there exists $\kappa > 0$ satisfying $D - d \ge \kappa H$.

Since $\log \epsilon < 0$, it follows that
\[
\epsilon^{D-d}
=
\exp\!\big((D-d)\log \epsilon\big)
\le
\exp\!\big(-\kappa H |\log \epsilon|\big).
\]

Absorbing constant and subexponential terms into the exponent,
there exists $c>0$ such that
\[
\frac{\mu(\mathcal{N}_\epsilon(\mathcal{M}_{\mathrm{traj}}))}{\mu(\mathcal{X})}
\le \exp(-cH),
\]
which completes the proof.
\end{proof}

\paragraph{Discussion.}
Lemma~\ref{lem:vanishing_mass} states that,
under uniform exploration of the trajectory space,
the probability of encountering even an approximately feasible trajectory
decays exponentially with the planning horizon.
This phenomenon arises from the mismatch between the rapidly growing ambient dimension
and the low intrinsic dimensionality of feasible behaviors,
and captures a fundamental scaling challenge faced by long-horizon
action-space planning.

